% Intracellular depletion --- plan \S3.3, "\S3".
% Sources: 01_opening.tex:137 (framing question, moved here), :58;
% 02_replicator_suppressor.tex:50-70, :76-92, :99-103, :107-154 (TikZ 1),
% :213-259 (TikZ fate regions), :263-366 (TikZ level diagonals), :370-393,
% :493-538.  The duplicate TikZ state space at 02:158-205 is deleted; the
% printed 15x15 matrix at 02:397-489 is replaced by prose plus the level
% diagram; the listing at 02:542-677 goes to Appendix B.
% Renames applied: alpha -> \varsigma, Y_t -> Q_t, theta_{i,j} -> \pi_{i,j};
% the abstract rates alpha_i, beta_{i,j} are written out as \lambda i and
% \varsigma ij.

\section{Intracellular depletion: replicator and suppressor}
\label{path:sec:intracellular}

A number of replicative units found an infection within a host cell, and
their dynamics accord with a continuous-time birth process; what is the
probability that at time $t$ at least one unit is still alive, so that the
infection has not died out? For a linear process this depends only on the
criticality --- whether a given unit is expected to leave behind more or
less than a replacement of itself --- and \cref{m:sec:bd,m:sec:gw} settle
it. \Cref{m:app:smallpops} asks the question this section asks, in its
linear form: how much difference does the size of the founding group make?
What is sought here is that same non-extinction probability through time
when the dynamics are nonlinear, and with it the effect of the nonlinearity
on survivability at varied initial counts $X_0=n$.

\subsection{The granule store}
\label{path:sec:granules}

Most of the handful of leucocyte types known to science carry packets of
toxic material called granules; granulocytes not least. Their contents are
generally oxidative, and the immune cell employs them to attack invading
pathogens \cite{jones2012subversion,cowley2011immunity}. The leucocytes of
interest in this research, macrophages and neutrophils, both carry granules,
but they carry different numbers of them and use them differently.
Neutrophils are of the category granulocyte and, as the name suggests, hold
more granules than macrophages; they lyse themselves to release those
inflammatory granules into the extracellular medium. Macrophages instead
employ their granules intracellularly, merging them with the phagosomal
membranes that harbour pathogen after engulfing has taken place --- a
compartment some pathogens escape rather than survive
\cite{golovliov2003attenuated}.
% NEEDS-BIOLOGY: whether neutrophils also employ granules intracellularly,
% in addition to the lytic release route.  The source asserted this as a
% belief ("I believe Neutrophils dothis as well") with no citation, and it
% is not checkable from the thesis tree.
% Candidate sources already copied into references.bib: cowley2011immunity,
% jones2012subversion.

It is this dynamic which forms the basis of the model presented here.
Picture a leucocyte which has just engulfed some number of pathogen and
which holds, besides, some number of toxin-containing granules, assumed to
be depleted in the process of being used to kill. Because the granules are
finite in number and depletable, a larger incident cohort of replicators may
have a greater chance of eventual survival. This is the intracellular form
of the dose-dependence that stochastic models of \ft{} and of
\textit{Salmonella} have found at the level of the whole host \cite{carruthers2020stochastic,gillard2014modeling,wood2014dose,brown2006intracellular}.

\begin{figure}[htbp]
    \centering
    \begin{subfigure}{0.40\textwidth}
        \centering
        \includegraphics[width=\textwidth]{RS1.png}
        \caption{}
    \end{subfigure}
    \hspace{0.03\textwidth}
    \begin{subfigure}{0.40\textwidth}
        \centering
        \includegraphics[width=\textwidth]{RS2.png}
        \caption{}
    \end{subfigure}

    \caption{The two opposing forces, simulated. Intracellular granules are
    in red, pathogenic particles in green; when one of each makes contact,
    both are annihilated, and that mutual annihilation is the depletion the
    model turns on. (a) Early, with the store still nearly full. (b) Later
    in the same run, with most of the store spent: the pathogen count is
    little changed, but the suppression rate $\varsigma i j$ has fallen with
    $j$, so what remains of the defence can no longer keep pace with the
    birth rate.}
    \label{path:fig:rs}
\end{figure}

\Cref{path:fig:rs} shows the two forces in a simulation. The point to take
from it is that the granules are consumed as they act: the red population in
panel~(b) is smaller than in panel~(a), and it is that reduction, not any
change in the pathogen, which lets the green population run away.

\subsection{The replicator--suppressor process}
\label{path:sec:rsmodel}

We explore a simple model representing the interplay between the two
opposing forces: replicative pathogenic particles fighting for survival
against their host cell's inbuilt defence mechanism. The model constitutes
two unit types, and from a defined initial condition the dynamic between
them results in either pathogen extinction or the overthrowing of immunity,
which here appears as unbounded growth in the replicator count.

\Needspace{16\baselineskip}
\begin{definition}[Replicator--suppressor process]
\label{path:def:rs}
Let $X_t$ count the infectious agents inside one host cell and $Q_t$ the
remaining number of within-cell granules. The process
$\mathcal{X}=\{(X_t,Q_t):t\ge0\}$ is the continuous-time Markov chain on
\begin{equation}
\mathcal{S}=\bigl\{(i,j)\;:\;i\in\mathbb{N}\cup\{0\},\;j\in\{0,\dots,N\}\bigr\}
\label{path:eq:statespace}
\end{equation}
with, for $(i,j)\in\mathcal{S}$,
\begin{equation}
\begin{aligned}
\pr{\Delta X_t=1,\;\Delta Q_t=0 \mid X_t=i,\,Q_t=j}
  &= \lambda i\,\Delta t+o(\Delta t),\\
\pr{\Delta X_t=-1,\;\Delta Q_t=-1 \mid X_t=i,\,Q_t=j}
  &= \varsigma ij\,\Delta t+o(\Delta t),\\
\pr{\Delta X_t=0,\;\Delta Q_t=0 \mid X_t=i,\,Q_t=j}
  &= 1-\lrb{\lambda i+\varsigma ij}\Delta t+o(\Delta t).
\end{aligned}
\label{path:eq:rsrates}
\end{equation}
\end{definition}

\begin{figure}[!h]
\centering
\begin{tikzpicture}[
  scale=0.85, >=Stealth, font=\footnotesize,
  st/.style={circle, draw, thick, fill=white,
             minimum size=4.4mm, inner sep=0pt, outer sep=0.4pt},
  abx/.style={rectangle, rounded corners=1.1pt, draw=Rup2, thick,
              fill=Rup2, minimum size=4.0mm, inner sep=0pt, outer sep=0.4pt},
  abq/.style={circle, draw=Rec2, thick, double, double distance=0.7pt,
              fill=Rec2fill, minimum size=4.4mm, inner sep=0pt, outer sep=0.4pt},
]
\draw[thick, color=inksoft] (0,0)--(7.7,0);
\draw[thick, color=inksoft] (0,0)--(0,4.7);
\node[below] at (3.85,-1.08) {$X_t$ \textit{replicators}};
\node[rotate=90] at (-1.18,2.35) {$Q_t$ \textit{granules}};
\foreach \x in {0,1,2,3,4,5,6,7}
  \node[below=4.0mm, inner sep=0pt, text=inksoft] at (\x,0) {\x};
\foreach \y in {1,2,3,4}
  \node[left=3.6mm, inner sep=0pt, text=inksoft] at (0,\y) {\y};

\foreach \y in {0,1,2,3,4}
  \node[abx] (s-0-\y) at (0,\y) {};
\foreach \x in {1,2,3,4,5,6,7}
  \node[abq] (s-\x-0) at (\x,0) {};
\foreach \x in {1,2,3,4,5,6,7}
  \foreach \y in {1,2,3,4}
    \node[st] (s-\x-\y) at (\x,\y) {};

\foreach \x/\xn in {1/2,2/3,3/4,4/5,5/6,6/7} {
  \foreach \y in {0,1,2,3,4}
    \draw[->, thick, na1] (s-\x-\y) -- (s-\xn-\y);
}
\foreach \y in {0,1,2,3,4}
  \draw[->, thick, na1, opacity=0.25] (s-7-\y.east) -- ++(0.48,0);

\foreach \x/\px in {1/0,2/1,3/2,4/3,5/4,6/5,7/6} {
  \foreach \y/\py in {1/0,2/1,3/2,4/3}
    \draw[->, thick, mut1] (s-\x-\y) -- (s-\px-\py);
}

\node[font=\scriptsize, text=na1, inner sep=0.6pt, fill=white]
  at (4.5,3.28) {$\lambda i$};
\node[font=\scriptsize, text=mut1, inner sep=0.6pt, fill=white]
  at (3.22,2.68) {$\varsigma ij$};

\begin{scope}[shift={(8.45,3.45)}, font=\scriptsize]
\draw[->, thick, na1] (0,0.55) -- (0.7,0.55);
\node[right, text=inksoft] at (0.8,0.55) {birth};
\draw[->, thick, mut1] (0.15,0.15) -- (-0.15,-0.15);
\node[right, text=inksoft] at (0.8,0.0) {suppression};
\node[abx] at (0.35,-0.58) {};
\node[right, text=inksoft] at (0.8,-0.58) {$X_t=0$, extinct};
\node[abq] at (0.35,-1.18) {};
\node[right, text=inksoft] at (0.8,-1.18) {$Q_t=0$, store spent};
\end{scope}
\end{tikzpicture}
\caption{The state space $\mathcal{S}$ of \cref{path:def:rs}. Births
move right and spend no granule; suppressions move down and left, spending
one of each. The column $X_t=0$ is extinction; the row $Q_t=0$ is a spent
store, from which only births remain, so survival is certain. That is why
the diagonal $i=j$ matters.}
\label{path:fig:statespace}
\end{figure}
Units of $X_t$ divide and increase according to the standard mechanism of a
birth process. Suppressive interactions occur at a rate proportional to the
product of the two counts and result in a single elimination of each type;
it is that product which makes the process nonlinear, and it is the only
nonlinearity in the model. Because each suppression costs one of each, a
cohort that outnumbers the store cannot be cleared.

In reality growth will not be truly without limit --- intracellular capacity
will eventually impose a threshold on further expansion without cell
bursting --- but the model is not run beyond the point at which the granules
are exhausted, so the idealisation costs nothing that is used below.
\Cref{path:fig:statespace} is the chain.

Our interest is in the relationship between incident pathogen count and
infection survivability, that is in
\begin{equation}
\pr{\text{infection survives}} = 1-\lim_{t\to\infty}\pr{X_t=0},
\label{path:eq:survival}
\end{equation}
which is the probability of late-time non-extinction. Survival is assured on
any path that exhausts the antagonistic granules: once $Q_t=0$ the process is
a pure birth process, and extinction is impossible thereafter.

\begin{remark}[No replenishment]
\label{path:rem:noreplenish}
The granule store is fixed at $Q_0$ and never supplemented. One may devise a
model in which granules are produced at some rate by the cell, and for a
long infection that would be the right model; the version analysed here
takes the initial immune capacity as given, which is the regime in which a
single engulfment episode is decided. Every result below is a statement
about that regime, and the certainty threshold of
\cref{path:prop:certainty} in particular depends on it.
\end{remark}

The asymmetry of the two arrow families is the whole model: a birth spends
nothing, whereas a suppression spends a granule, so the store can only fall
and each unit of it can be used once.

\begin{figure}[htbp]
\centering
\begin{tikzpicture}[
  scale=0.58, font=\footnotesize,
  st/.style={circle, draw, thick, fill=white,
             minimum size=3.8mm, inner sep=0pt, outer sep=0.3pt},
  abx/.style={rectangle, rounded corners=0.9pt, draw=Rup2, thick,
              fill=Rup2, minimum size=3.4mm, inner sep=0pt, outer sep=0.3pt},
  abq/.style={circle, draw=Rec2, thick, double, double distance=0.65pt,
              fill=Rec2fill, minimum size=3.8mm, inner sep=0pt, outer sep=0.3pt},
]
\draw[thick, color=inksoft] (0,0)--(9.4,0);
\draw[thick, color=inksoft] (0,0)--(0,6.5);
\node[below] at (4.5,-1.05) {$X_t$ \textit{replicators}};
\node[rotate=90] at (-1.22,3.2) {$Q_t$ \textit{granules}};
\foreach \x in {0,1,2,3,4,5,6,7,8}
  \node[below=3.6mm, inner sep=0pt, text=inksoft] at (\x,0) {\x};
\node[below=3.6mm, inner sep=0pt, text=inksoft] at (9,0) {$\cdots$};
\foreach \y in {1,2,3,4,5,6}
  \node[left=3.4mm, inner sep=0pt, text=inksoft] at (0,\y) {\y};

\begin{scope}
\clip (0,0) rectangle (9.35,6.5);
\fill[Rec2fill] (0,0) -- (6.5,6.5) -- (9.35,6.5) -- (9.35,0) -- cycle;
\fill[nb1, opacity=0.45] (0.42,6.5) -- (6.5,6.5) -- (0.42,0.42) -- cycle;
\fill[Rup2fill] (0,0) rectangle (0.38,6.5);
\end{scope}

\foreach \y in {0,1,2,3,4,5,6}
  \node[abx] at (0,\y) {};
\foreach \x in {1,2,3,4,5,6,7,8}
  \node[abq] at (\x,0) {};
\foreach \x in {1,2,3,4,5,6,7,8}
  \foreach \y in {1,2,3,4,5,6}
    \node[st] at (\x,\y) {};
\foreach \y in {1,2,3,4,5,6}
  \node at (9,\y) {$\cdots$};
\node[abq] at (9,0) {};

\begin{scope}[shift={(10.15,4.7)}, font=\scriptsize]
\fill[Rec2fill] (0,-0.12) rectangle (0.36,0.16);
\draw[Rec2, thin] (0,-0.12) rectangle (0.36,0.16);
\node[anchor=west, text=inksoft] at (0.48,0.02) {certain, $i>j$};
\fill[nb1, opacity=0.45] (0,-0.67) rectangle (0.36,-0.39);
\draw[inksoft, thin] (0,-0.67) rectangle (0.36,-0.39);
\node[anchor=west, text=inksoft] at (0.48,-0.53) {contested, $1\le i\le j$};
\node[abx] at (0.18,-1.12) {};
\node[anchor=west, text=inksoft] at (0.48,-1.12) {extinct, $i=0$};
\end{scope}
\end{tikzpicture}
\caption{The three fates, read off \cref{path:fig:statespace}. In $i>j$
survival is certain: extinction would need $i$ suppressions and only $j$
are available (\cref{path:prop:certainty}). In $1\le i\le j$ both outcomes
have positive probability, and $\pi_{i,j}$ is what the rest of this section
computes. The boundary is the diagonal, and it is sharp.}
\label{path:fig:fate}
\end{figure}

\Cref{path:fig:fate} is the section's headline. A cohort that outnumbers the
store cannot be cleared, however unfavourable the rate constants; a cohort
that does not outnumber it may still win, and with what probability is the
question.

\subsection{First-step analysis}
\label{path:sec:rsfirststep}

The process is defined by its one-step transition probabilities,
\[
p_{(i,j)\to(k,l)}(\Delta t)
= \pr{(X_{t+\Delta t},Q_{t+\Delta t})=(k,l)\mid (X_t,Q_t)=(i,j)},
\]
and enumerating the steps available from an arbitrary state $(i,j)$ gives
\begin{equation}
p_{(i,j)\to(k,l)}(\Delta t)=
\begin{cases}
\lambda i\,\Delta t + o(\Delta t), & (k,l)=(i+1,j),\\
\varsigma ij\,\Delta t + o(\Delta t), & (k,l)=(i-1,j-1),\\
1-\lrb{\lambda i+\varsigma ij}\Delta t + o(\Delta t), & (k,l)=(i,j),
\end{cases}
\label{path:eq:onestep}
\end{equation}
as $\Delta t\to0$. Write $\pi_{i,j}$ for the probability of fixation in the
set $\{X_t=0\}$ from the state $(i,j)$. Conditioning on the first jump, and
using the first-step principle of \cref{m:prop:firststep} in the
continuous-time form \cref{m:eq:firststep},
\begin{equation}
\pi_{i,j} = \hat\alpha_j\,\pi_{i+1,j} + \hat\beta_j\,\pi_{i-1,j-1},
\qquad
\hat\alpha_j=\frac{\lambda}{\lambda+\varsigma j},
\qquad
\hat\beta_j=\frac{\varsigma j}{\lambda+\varsigma j},
\label{path:eq:firststep}
\end{equation}
with $\pi_{0,j}=1$ for every $j$ and $\pi_{i,0}=0$ for every $i\ge1$. The
jump probabilities depend on $j$ alone: the competing rates are $\lambda i$
and $\varsigma ij$, and the common factor $i$ cancels. The replicator count
therefore affects the outcome only through how many suppressions are needed,
never through how likely the next event is to be one.

Collecting \cref{path:eq:firststep} over all transient states gives the
linear system
\begin{equation}
\mathbf{p}^{(0)} = \mathbf{A}\,\mathbf{p}^{(0)} + \mathbf{c},
\label{path:eq:matrixsystem}
\end{equation}
in which $\mathbf{p}^{(0)}$ collects the unknowns $\pi_{i,j}$ for
$1\le i\le j\le N$, the matrix $\mathbf{A}$ carries the coefficients
$\hat\alpha_j$ and $\hat\beta_j$, and the vector $\mathbf{c}$ carries the
single inhomogeneous term $\hat\beta_j$ contributed by $\pi_{0,j-1}=1$ at
$i=1$. The system is block lower-triangular when the states are ordered by
the level $j-i$, because the two states appearing on the right of
\cref{path:eq:firststep} sit one at level $j-i-1$ and one at level $j-i$;
those levels are drawn in \cref{path:fig:levels}. Nothing beyond that
structure is needed: it makes the system solvable level by level, and the
solver of \cref{path:app:computational} exploits nothing else.

\begin{figure}[htbp]
\centering
\begin{tikzpicture}[
  scale=0.68, font=\footnotesize,
  st/.style={circle, draw, thick, fill=white,
             minimum size=4.0mm, inner sep=0pt, outer sep=0.3pt},
  abx/.style={rectangle, rounded corners=0.9pt, draw=Rup2, thick,
              fill=Rup2, minimum size=3.6mm, inner sep=0pt, outer sep=0.3pt},
  abq/.style={circle, draw=Rec2, thick, double, double distance=0.65pt,
              fill=Rec2fill, minimum size=4.0mm, inner sep=0pt, outer sep=0.3pt},
]
\draw[thick, color=inksoft] (0,0)--(8.2,0);
\draw[thick, color=inksoft] (0,0)--(0,6.5);
\node[below] at (3.7,-1.05) {$X_t$ \textit{replicators}};
\node[rotate=90] at (-1.18,3.1) {$Q_t$ \textit{granules}};
\foreach \x in {0,1,2,3,4,5,6,7}
  \node[below=3.6mm, inner sep=0pt, text=inksoft] at (\x,0) {\x};
\foreach \y in {1,2,3,4,5,6}
  \node[left=3.4mm, inner sep=0pt, text=inksoft] at (0,\y) {\y};

\begin{scope}
\clip (0,0) rectangle (8.15,6.5);
\foreach \k in {0,1,2,3,4,5} {
  \draw[na1, thick, opacity=0.50]
    (0.5,{0.5+\k}) -- ({6.5-\k},6.5);
}
\end{scope}

\foreach \y in {0,1,2,3,4,5,6}
  \node[abx] at (0,\y) {};
\foreach \x in {1,2,3,4,5,6,7}
  \node[abq] at (\x,0) {};
\foreach \x in {1,2,3,4,5,6,7}
  \foreach \y in {1,2,3,4,5,6}
    \node[st] at (\x,\y) {};
\foreach \y in {1,2,3,4,5,6}
  \node at (8,\y) {$\cdots$};

\begin{scope}[shift={(8.55,5.55)}, font=\scriptsize]
\node[anchor=west, text=inksoft] at (0,0.15) {level $j-i$};
\foreach \k in {0,1,2,3,4,5} {
  \draw[na1, thick] (0,{-0.38-0.42*\k}) -- (0.48,{-0.38-0.42*\k});
  \node[anchor=west, text=na1] at (0.58,{-0.38-0.42*\k}) {\k};
}
\end{scope}
\end{tikzpicture}
\caption{The level structure of the first-step system
\cref{path:eq:firststep}. States are grouped by the level $j-i$. A
suppression stays on its diagonal; a birth moves one level down. No
transition raises the level, so $\mathbf{A}$ in \cref{path:eq:matrixsystem}
is block lower-triangular in this ordering and the levels may be solved in
turn, starting from level~$0$.}
\label{path:fig:levels}
\end{figure}

Level~0 is the diagonal $i=j$, and it closes on itself. That is enough for
two results.

\begin{proposition}[Certainty threshold]
\label{path:prop:certainty}
$\pi_{i,j}=0$ whenever $i>j$: the infection survives with probability one as
soon as the incident cohort exceeds the granule store. For $1\le i\le j$ we
have $\pi_{i,j}\in(0,1)$.
\end{proposition}

\begin{proof}
Each suppression removes one pathogen and one granule; a birth removes no
granule. Reaching $X=0$ from $(i,j)$ therefore requires at least $i$
suppressions, so $j\ge i$ as claimed.

Now take $1\le i\le j$. A run of $i$ suppressions reaches $(0,j-i)$ and has
positive probability, since $\hat\beta_j>0$ at every $j\ge1$; thus
$\pi_{i,j}>0$. A run of $j-i+1$ births reaches $(j+1,j)$, where
$\pi=0$ by the first claim, and has positive probability since
$\hat\alpha_j>0$; thus $\pi_{i,j}<1$.
\end{proof}

The boundary is the diagonal drawn in \cref{path:fig:fate}, and nothing
about the rate constants moves it. Consequently the only remaining question
on the contested side is how $\pi_{i,j}$ varies with dose. Monotonicity in
$i$ and in $j$ was checked numerically over
$\varsigma\in\{0.1,0.2,0.3\}$ at $N=20$, and both hold everywhere
(\cref{path:app:verification}).

\begin{proposition}[The diagonal law]
\label{path:prop:diagonal}
Write $\nu=\lambda/\varsigma$. On the matched diagonal,
\begin{equation}
\pi_{j,j}
= \prod_{m=1}^{j}\frac{\varsigma m}{\lambda+\varsigma m}
= \frac{j!\,\Gamma(\nu+1)}{\Gamma(j+\nu+1)}
= \binom{j+\nu}{j}^{-1}
\;\sim\; \Gamma(\nu+1)\,j^{-\nu}
\qquad (j\to\infty),
\label{path:eq:diagonal}
\end{equation}
so a cohort matched to its own granule store fails with a probability
falling only polynomially in cohort size, at exponent $\lambda/\varsigma$.
When $\lambda=\varsigma$ this is exactly $1/(j+1)$.
\end{proposition}

\begin{proof}
From $(j,j)$ a birth reaches $(j+1,j)$, where $\pi=0$ by
\cref{path:prop:certainty}. Hence $\pi_{j,j}=\hat\beta_j\,\pi_{j-1,j-1}$
with $\hat\beta_j=\varsigma j/(\lambda+\varsigma j)$, and $\pi_{0,0}=1$.
Iterating gives the product, and the ratio of Gamma functions follows from
$\prod_{m=1}^{j}\varsigma m=\varsigma^{j}j!$ and
$\prod_{m=1}^{j}(\lambda+\varsigma m)=\varsigma^{j}\Gamma(j+\nu+1)/\Gamma(\nu+1)$.
The asymptotic is the standard ratio limit
$\Gamma(j+1)/\Gamma(j+\nu+1)\sim j^{-\nu}$.
\end{proof}

\begin{table}[tbp]
\centering
\small
\begin{tabular}{@{}lcccc@{}}
\toprule
$\varsigma$ & $\nu=\lambda/\varsigma$ & $\pi_{5,5}$ & $\pi_{10,10}$ & $\pi_{25,25}$ \\
\midrule
$0.1$ & $10$   & $3.33\times10^{-4}$ & $5.41\times10^{-6}$ & $5.45\times10^{-9}$ \\
$0.2$ & $5$    & $3.97\times10^{-3}$ & $3.33\times10^{-4}$ & $7.02\times10^{-6}$ \\
$0.3$ & $10/3$ & $1.34\times10^{-2}$ & $2.26\times10^{-3}$ & $1.54\times10^{-4}$ \\
$1.0$ & $1$    & $1.67\times10^{-1}$ & $9.09\times10^{-2}$ & $3.85\times10^{-2}$ \\
\bottomrule
\end{tabular}
\caption{The diagonal law \cref{path:eq:diagonal} at $\lambda=1$. Reading
across a row shows the decay in cohort size at fixed rates; reading down a
column shows how much slower that decay becomes as suppression strengthens
relative to birth. The bottom row is the borderline case $\lambda=\varsigma$,
where the failure probability is exactly $1/(j+1)$ --- a cohort of
twenty-five matched to twenty-five granules still fails one time in
twenty-six.}
\label{path:tab:diagonal}
\end{table}

\Cref{path:tab:diagonal} shows why \cref{path:eq:diagonal} is worth stating
separately. Doubling the cohort on the diagonal does not halve the failure
probability: it divides it by $2^{\nu}$, and $\nu=\lambda/\varsigma$ is a
property of the cell rather than of the dose. Safety in numbers is therefore
real but slow, and qualitatively weaker than the certainty that arrives the
moment the cohort crosses the store.

\subsection{The survival maps}
\label{path:sec:rsnumerics}

Solving \cref{path:eq:matrixsystem} over the whole grid gives the survival
probability $1-\pi_{i,j}$ for every initial condition.

\begin{figure}[htbp]
    \centering
    \includegraphics[width=\textwidth]{repsupp_N20.pdf}
    \caption{Survival probability $1-\pi_{i,j}$ at $\lambda=1$ and $N=20$,
    for $\varsigma=0.1$, $0.2$ and $0.3$ (left to right). Each lower panel
    is the full grid of initial conditions, with the dashed line marking the
    diagonal $i=j$; each upper panel is the marginal at a full store,
    $1-\pi_{i,N}$, which is the case a freshly engulfing macrophage
    presents. Everything below the diagonal is certain survival
    (\cref{path:prop:certainty}); the transition above it steepens as
    $\varsigma$ falls, so that a weakly armed cell is not merely less
    effective but less discriminating about dose.}
    \label{path:fig:repsupp20}
\end{figure}

\begin{figure}[htbp]
    \centering
    \includegraphics[width=0.92\textwidth]{repsupp_N10.pdf}
    \caption{The same quantity at $N=10$, again with $\lambda=1$ and
    $\varsigma=0.1$, $0.2$, $0.3$. Halving the store does not halve the dose
    a pathogen needs: the certainty boundary moves with $N$, so the
    requirement is a matter of the ratio of cohort to store rather than of
    the cohort alone.}
    \label{path:fig:repsupp10}
\end{figure}

\Cref{path:fig:repsupp20,path:fig:repsupp10} show the maps at $N=20$ and
$N=10$. Two features carry the section. The first is the flat region of certain survival
below the diagonal, which is \cref{path:prop:certainty} drawn: it is not
that survival becomes likely there, but that it becomes certain. The second
is the marginal bar chart above each map, the case $j=N$ of a full store, in
which the rise from near-zero to
near-one is spread over several units of dose rather than concentrated at
one --- the polynomial decay of \cref{path:eq:diagonal}, seen from the side.

For a pathogen, the reading is direct. If a host cell can be entered by
several agents at once rather than one at a time, and the store is finite,
then arriving together is worth more than arriving in succession by exactly
the margin these maps display. That is the intracellular version of the
flooding advantage, and \cref{path:sec:removal} builds its extracellular
counterpart.

\FloatBarrier
